\section{Dawson's chess on the cylinder}\label{sec:dawson}

Dawson's chess \citep{Dawson1935} is played on a $3\times n$ board:
\begin{itemize}
\item White pawns fill rank~1, Black pawns fill rank~3, and rank~2 is empty;
\item pawns move as in chess, and captures are compulsory;
\item a player who cannot move loses (normal play).
\end{itemize}
\citet{GuySmith1956} identified normal play with the octal game $0.137$. Its Grundy values $g(n)$ have
period $34$, with last exception $n=51$ (classical; reproduced for $n\le300$, \status{checked-2}).

\paragraph{Correction.} The 2025 paper's \S8 is wrong in three ways (\status{checked-2}):
\begin{itemize}
\item its move rule allows no first move;
\item its recurrence omits the end-file move and $g(1)=1$, and agrees with the true $g$ at only $31$ of the
  values $n=0,\dots,300$;
\item its cylinder predictions are therefore wrong for $56$ values of $n\le300$.
\end{itemize}

\begin{theorem}[ring rule; \status{proved}, machine-checked for $n\le22$]\label{thm:ring}
On the full $3\times n$ cylinder with $n\ge3$, every first move forces exactly five plies of captures on the
pushed file~$f$. All lines end in the same position: files $f\pm1$ empty, a blocked pair on $f$, and a full flat
strip of $n-3$ files with the opponent to move. Hence the ring's Grundy value is $\operatorname{mex}\{g(n-3)\}$,
and the first player wins iff $g(n-3)=0$, that is for
$n=3,7,11,17,23,27,31,37,41,45,57,\dots$
\end{theorem}

\paragraph{Misère play.} Dawson posed the game in misère form, where the last player to move loses. That
version is still unsolved.
\begin{itemize}
\item An independent implementation of the Plambeck--Siegel quotient algorithm \citep{PlambeckSiegel2008}
  reproduces their partial quotient for strips of length $\le32$ exactly: order $638$, with $109$ P-positions.
  It agrees with exhaustive search on $7.4$ million positions (\status{checked-2}).
\item The next quotient did not converge within our budget.
\item The partizan values that a direct solver produces for each board are not additive (for example $D_5=\star$
  while $g(5)=3$). The additive invariant is the impartial macro-move value.
\end{itemize}
